\subsection{The constant $\kappa$ of the reach: how it is determined}
\begin{problem}[title={What $\kappa$ is}]
$\kappa$ is the relative conductivity $K_r=K/K_s$ at which the crop can no longer dry the soil at the sensor within a season. It is the only constant of \cref{eq:s-hc}. It is \emph{not} measured directly: it is found from the simulated triggers, in five steps.
\end{problem}
\begin{derivation}[title={How $\kappa$ was determined, step by step}]
\textbf{Step 1: the experiment.} Experiment E2 (197 valid runs) sets the trigger to 20 values from 50 to 10\,000\,cm on each of the ten soils. Everything else is fixed: the sensor at 20\,cm, a dose of 2.4\,mm per event, surface application, the crop file \texttt{crops\_conv\_N100.dai} (the cereal rotation of the ensemble under conventional tillage with 100\,kg\,N\,ha$^{-1}$), and 20 years of Taastrup weather (1980--1999).

\textbf{Step 2: when does a trigger fire?} A trigger is counted as firing on a soil if it brings at least 5\,mm of irrigation a year (20-year mean).

\textbf{Step 3: the observed reach of each soil.} For each soil the reach lies between the driest trigger that still fires ($h_{\mathrm{lo}}$) and the next drier one, which does not ($h_{\mathrm{hi}}$); if even 10\,000\,cm fires, only $h_{\mathrm{lo}}$ is known. Sand: 100--150\,cm; loamy sand: 200--250\,cm; sandy loam: 800--1\,000\,cm; loam and sandy clay loam: 5\,000--10\,000\,cm; the five finer or coarse-loamy soils: above 10\,000\,cm.

\textbf{Step 4: one $\kappa$ for all soils.} For a trial $\kappa$, \cref{eq:s-hc} gives a reach for every soil. Its distance to the observed interval is measured in decades (zero if inside, the distance to the nearer end if outside), squared and added over the ten soils.

\textbf{Step 5: the best $\kappa$.} The sum is minimised over $\log_{10}\kappa$ with the Nelder--Mead method. The minimum is at $\kappa=3.47\times10^{-10}$.
\end{derivation}
\begin{keyresult}[title={How certain $\kappa$ is}]
\centering\footnotesize
\begin{tabular}{>{\raggedright\arraybackslash}p{4.3cm}>{\raggedright\arraybackslash}p{5.2cm}>{\raggedright\arraybackslash}p{4.0cm}}
\toprule
Test & Result & What it means\\\midrule
Bootstrap over soils (1\,000 resamples) & $\kappa$ between $1.0\times10^{-10}$ and $2.8\times10^{-9}$ (95\,\%) & $\kappa$ is known to about one decade\\
Leaving each soil out & $2.3\times10^{-10}$ to $7.3\times10^{-10}$ & no single soil decides $\kappa$\\
Firing threshold 1 or 2\,mm & $2.7\times10^{-10}$ & a stricter count changes $\kappa$ little\\
Firing threshold 5\,mm (used) & $3.5\times10^{-10}$ & \\
Firing threshold 10\,mm & $7.3\times10^{-10}$ & \\
Firing threshold 20\,mm & $1.3\times10^{-9}$ & \\
Firing threshold 40\,mm & $2.5\times10^{-9}$ & a looser count raises $\kappa$\\\bottomrule
\end{tabular}
\end{keyresult}
\begin{pins}[title={How much an error in $\kappa$ moves the reach}]
\pin{1}{Because $\hc\propto\kappa^{-1/E}$, a tenfold error in $\kappa$ changes the reach by the factor $10^{1/E}$: 1.45 in sand, 1.56 in loamy sand, 1.73 in sandy loam, 1.97 in loam and 2.45 in silty clay loam. The reach is least sensitive exactly in the coarse soils where it matters most (where it decides whether a 300\,cm trigger works).}
\pin{2}{Moving the firing threshold from 5 to 20\,mm a year raises $\kappa$ about fourfold and shortens the reaches by a factor of 1.2 (sand) to 1.4 (silty clay loam); sand and loamy sand remain blind to a 300\,cm trigger with any threshold tested.}
\end{pins}
\begin{keyresult}[title={What $\kappa$ depends on}]
\textbf{Not the soil.} One value fits ten soils from sand to silty clay loam within a factor of 2.3, and the soil enters only through $\alpha$, $n$ and $\ell$.

\textbf{The crop: yes, in principle.} $\kappa$ expresses when the roots can no longer extract water from the soil around the sensor. It depends on the root length density at the sensor depth, the potential transpiration of the crop and the lowest root water potential the crop can sustain. A crop with denser or deeper roots, or one that tolerates lower water potentials, can dry the soil further, which means a smaller $\kappa$ and a longer reach. All runs of E2 used one cereal rotation, so the size of this effect was not measured.

\textbf{The climate: yes.} A higher evaporative demand stresses the crop earlier; the weather was one series (Taastrup), so this was not measured either.

\textbf{The sensor depth and the dose: yes.} The sensor was at 20\,cm, where the roots are dense; a deeper sensor sees fewer roots and dries less (in E3 a deeper sensor did not make sand irrigate at 300\,cm). The dose was 2.4\,mm.

\textbf{In practice.} $\kappa$ is a constant of the crop--climate--sensor system, not of the soil. For another crop or climate it is determined in the same way, with a trigger experiment of the kind of E2 on two or three soils, which the Daisy workbench of SWAFL Desktop can run.
\end{keyresult}

\subsection{Every other constant of the article}
{\small\begin{longtable}{>{\raggedright\arraybackslash}p{\dimexpr 0.1711\linewidth-2\tabcolsep\relax}>{\raggedright\arraybackslash}p{\dimexpr 0.1513\linewidth-2\tabcolsep\relax}>{\raggedright\arraybackslash}p{\dimexpr 0.3553\linewidth-2\tabcolsep\relax}>{\raggedright\arraybackslash}p{\dimexpr 0.3224\linewidth-2\tabcolsep\relax}}
\toprule Constant & Value & How it was obtained & What it depends on; sensitivity\\\midrule\endhead
$\ell$, Mualem connectivity & 0.5 & conventional value of \citet{Mualem1976}, not fitted & pore structure; a larger $\ell$ steepens the dry-end conductivity and shortens the reach\\
Firing threshold & 5\,mm\,yr$^{-1}$ & chosen: the smallest amount clearly above the numerical noise of the 20-year means & see the table above: $\kappa$ changes 2.7 to $25\times10^{-10}$ for 1 to 40\,mm\\
Factor $\tfrac12$ of the blind rule & 0.5 & read from the data: the driest trigger that still fires lies at $\htr/\hc=0.47$ (loamy sand), 0.49 (sandy loam), 0.50 (sand), 0.54 (loam) and 0.66 (sandy clay loam) & a margin, not a fitted constant; below one half every soil still irrigates\\
$\alpha$, $n$ of each soil & Table~A1 & effective (profile-averaged) values of the soils of the ensemble & the soil; see the terms table\\
$K_s$ (only to test the absolute-conductivity law) & 1.68 to 712.8\,cm\,d$^{-1}$ & class means of \citet{Carsel1988} for eight soils & not used by \cref{eq:s-hc}\\
$\eta_\infty$ of the dose law & 0.10 to 0.45 & nonlinear least squares per soil over ten doses (standard errors 0.005 to 0.03) & the soil: how fast it drains once wetted\\
$d_D$ of the dose law & 3.2 to 5.2\,mm & same fit (standard errors 0.5 to 1.8\,mm; sandy loam least certain) & the crop and the atmosphere (one day of water use), not the soil ($\rho=0.07$ with $n$)\\
Rainfed baseline $X_0$ & per soil & experiment E0, same soil and weather without irrigation & the weather and the crop\\
Validity of the dry-end form & $\alpha\hc\ge29$ & checked for every soil: the neglected terms are below $(\alpha\hc)^{-n}<10^{-3}$ & holds for all ten soils\\\bottomrule
\end{longtable}}
